% Direct proof of the P06.1 analytic algebraic-split toric-form extension.
% The field class is exactly the one in GJR26, Corollary 5.12; this neither
% manufactures a splitting chart nor applies to a general analytic extension.

\begin{proof}[Proof of
\Cref{cor:p061-algebraic-split-analytic-toric-forms}]
We prove the product assertion; the assertion without a parameter is the
case where \(Y\) is a point.  Let \(\Omega\) be an ordinary open analytic
subspace of \(X\times_KY\), let \(h\in\operatorname{PH}(\Omega)\), and fix
\((x,y)\in\Omega\).  Choose a displayed ordinary neighbourhood \(U\) of
\(x\), its complete algebraic splitting extension \(L/K\), and the
displayed analytic identification of \(U_L\) with an ordinary open analytic
subspace of \(\mathcal X^{\mathrm{an}}\).  Put
\[
 \Omega_U=\Omega\cap(U\times_KY).
\]
After scalar extension, \((\Omega_U)_L\) is an ordinary open analytic
subspace of
\[
 \mathcal X^{\mathrm{an}}\times_LY_L.
\]
The space \(Y_L\) is good, strict, and quasi-smooth.  By scalar-extension
stability of classical pluriharmonicity
\cite[Proposition~7.2(6)]{GR25Classical}, the pullback \(h_L\) is
classically pluriharmonic on \((\Omega_U)_L\).  Hence
\Cref{cor:p061-split-toric-quasismooth-products} gives
\[
 h_L\in\operatorname{PL}_{\mathbf R}((\Omega_U)_L).
\]

The ordinary open space \(\Omega_U\) is good and strict, while \(L\) is
the completion of an algebraic extension of \(K\).  Thus ground-field
descent for \(\mathbf R\)-PL functions, with coefficient group
\(\mathbf R\), applies by \cite[Corollary~5.12]{GJR26}; it gives
\(h|_{\Omega_U}\in\operatorname{PL}_{\mathbf R}(\Omega_U)\).  Since
\((x,y)\) was arbitrary, \(h\) is locally \(\mathbf R\)-PL on \(\Omega\),
and hence \(\mathbf R\)-PL there.  Both \(h\) and \(-h\) are classically
psh.  By \cite[Theorem~4.3]{GR25Classical}, both signs are semipositive, so
\(h\) is Gubler--Jell--Rabinoff harmonic by \cite[Definition~7.2]{GJR26}.

Finally, transfer a classically pluriharmonic function through the displayed
analytic local isomorphism and apply the product assertion on the resulting
ordinary open subspace.  The conclusions transfer back through that local
isomorphism.
\end{proof}
